\documentclass[12pt,a4paper]{article}

\usepackage[utf8]{inputenc}
\usepackage[T1]{fontenc}
\usepackage[spanish,es-tabla]{babel}
\usepackage{geometry}
\usepackage{array}
\usepackage{longtable}
\usepackage{booktabs}
\usepackage{xcolor}
\usepackage{hyperref}
\usepackage{enumitem}
\usepackage{amsmath}

\geometry{left=2.4cm,right=2.4cm,top=2.4cm,bottom=2.4cm}
\setlength{\parindent}{0pt}
\setlength{\parskip}{6pt}
\setlength{\tabcolsep}{3pt}
\emergencystretch=2em
\setlist[itemize]{leftmargin=1.2cm}
\setlist[enumerate]{leftmargin=1.2cm}

\hypersetup{
  colorlinks=true,
  linkcolor=black,
  urlcolor=blue
}

\newcommand{\pendiente}{\textcolor{red}{Pendiente}}
\newcommand{\obs}{\textcolor{orange!70!black}{Observacion}}
\newcommand{\ok}{\textcolor{green!50!black}{Definido}}
\newcolumntype{L}[1]{>{\raggedright\arraybackslash}p{#1}}
\newcolumntype{R}[1]{>{\raggedleft\arraybackslash}p{#1}}

\begin{document}

\begin{center}
{\Large \textbf{Planificacion inicial del proyecto}}\\
{\large Capitulos I y II - Instalacion electrica domiciliaria}\\
{\normalsize Vivienda propia de 2 pisos}
\end{center}

\section*{1. Datos generales registrados}

\begin{longtable}{L{5.2cm}L{7.8cm}L{2.2cm}}
\toprule
\textbf{Dato} & \textbf{Respuesta} & \textbf{Estado} \\
\midrule
\endfirsthead
\toprule
\textbf{Dato} & \textbf{Respuesta} & \textbf{Estado} \\
\midrule
\endhead
Estudiante & Aquiles Taylor Ramos Yapo & \ok \\
Docente / ingeniero & Por completar & \pendiente \\
Fecha de avance & 01 de junio de 2026 & \ok \\
Direccion o referencia & Avenida Horacio con Jr. Marineros, manzana F7, lotes 11 y 12 & \ok \\
Distrito & San Miguel & \ok \\
Provincia & San Román & \ok \\
Departamento & Puno & \ok \\
Zona & Urbana & \ok \\
Tipo de vivienda real & Vivienda unifamiliar (casa grande) & \ok \\
Alcance academico & Vivienda unifamiliar para 7 miembros & \ok \\
Material predominante & Ladrillo y concreto & \ok \\
Area por piso & Primer piso con menor área construida que el segundo. Terreno: 134.18 m² & \obs \\
Empresa distribuidora & Electro Puno & \ok \\
Sistema asumido & Monofasico, 220 V, 60 Hz & \ok \\
\bottomrule
\end{longtable}

\section*{2. Criterio de alcance}

Aunque la vivienda real puede ser usada por mas de una familia, para fines academicos se delimitara el proyecto a una vivienda unifamiliar de dos pisos y 7 miembros. El trabajo se concentrara en la instalacion electrica interior: alumbrado, tomacorrientes, cocina, cargas de servicio, bomba de agua, tablero, protecciones y puesta a tierra propuesta.

No se tomara como definitivo ningun metrado ni longitud de circuito hasta contar con croquis y dimensiones aproximadas. Las areas quedan en observacion.

\section*{3. Distribucion arquitectonica preliminar}

\subsection*{3.1 Primer piso}

\begin{longtable}{L{4cm}R{2cm}R{3cm}L{5.2cm}}
\toprule
\textbf{Ambiente / zona} & \textbf{Focos} & \textbf{T.C.} & \textbf{Observacion} \\
\midrule
\endfirsthead
\toprule
\textbf{Ambiente / zona} & \textbf{Focos} & \textbf{T.C.} & \textbf{Observacion} \\
\midrule
\endhead
Cuarto 1 grande & 2 & 3 & Division de drywall sin tomacorriente \\
Cuarto 2 grande & 2 & 3 & Division de drywall sin tomacorriente \\
Cocina pequena & 2 & 1 & Licuadora y artefactos pequenos \\
Pasadizo & 3 & 0 & Alumbrado de circulacion \\
Escalera / plataforma intermedia & 1 & 0 & Foco entre primer y segundo piso \\
Bano & 0 & 0 & El bano esta en segundo piso \\
Bomba exterior & 0 & circuito dedicado & Bomba de agua en pozo exterior \\
\bottomrule
\end{longtable}

\subsection*{3.2 Segundo piso}

\begin{longtable}{L{4cm}R{2cm}L{3cm}L{5.2cm}}
\toprule
\textbf{Ambiente / zona} & \textbf{Focos} & \textbf{T.C. / carga} & \textbf{Observacion} \\
\midrule
\endfirsthead
\toprule
\textbf{Ambiente / zona} & \textbf{Focos} & \textbf{T.C. / carga} & \textbf{Observacion} \\
\midrule
\endhead
Sala & 2 & por definir & Computador y parlantes \\
Cuarto con cama 1 & 2 & por definir & Pantalla plana de 50 pulgadas \\
Cuarto con cama 2 & 2 & por definir & Pantalla plana de 50 pulgadas \\
Cocina amplia / zona de servicio & 2 & 4 & Microondas, lavadora, waflera y cargas de consumo mayor \\
Cuarto de uso varios & 2 & por definir & Actualmente en desuso \\
Bano & 1 & 1 & Tomacorriente protegido para usos varios \\
\bottomrule
\end{longtable}

\section*{4. Respuestas tecnicas relevantes}

\begin{longtable}{L{5cm}L{9.2cm}}
\toprule
\textbf{Tema} & \textbf{Respuesta registrada} \\
\midrule
\endfirsthead
\toprule
\textbf{Tema} & \textbf{Respuesta registrada} \\
\midrule
\endhead
Cocina del primer piso & Solo licuadora y artefactos pequenos o manejables. \\
Ducha electrica & No se considera. \\
Terma & Se usa terma solar. \\
Bomba de agua & Existe bomba exterior en el primer piso, asociada a pozo. \\
Lavadora & Prevista en el segundo piso. \\
Horno / microondas & Previsto en cocina amplia del segundo piso. \\
Waflera y equipos similares & Considerar como carga de consumo medio/alto en segundo piso. \\
Croquis & Pendiente; se dibujara en hoja y luego se digitalizara. \\
Dimensiones & Aproximadas, por no encontrarse actualmente en la vivienda. \\
Puesta a tierra actual & La vivienda no cuenta actualmente con puesta a tierra. \\
Puesta a tierra propuesta & Se recomienda proponerla desde el avance. \\
\bottomrule
\end{longtable}

\section*{5. Normativa base para sostener el avance}

\begin{longtable}{L{3.4cm}L{4.8cm}L{6cm}}
\toprule
\textbf{Norma} & \textbf{Referencia extraida} & \textbf{Uso en el proyecto} \\
\midrule
\endfirsthead
\toprule
\textbf{Norma} & \textbf{Referencia extraida} & \textbf{Uso en el proyecto} \\
\midrule
\endhead
RNE EM.010 & Articulos 1 y 2, paginas 1--2 & Sustenta que la instalacion electrica interior comprende acometida, alimentadores, tableros, circuitos derivados, sistemas de proteccion, medicion y puesta a tierra; aplica a viviendas. \\
CNE-U & Seccion 010, pagina 6 & Sustenta la obligatoriedad y alcance general del Codigo Nacional de Electricidad - Utilizacion. \\
CNE-U & Pagina 13 & Apoya la definicion de unidad de vivienda. \\
CNE-U & Seccion 050, pagina 62 & Apoya el criterio de demanda para viviendas unifamiliares cuando no se dispone de informacion completa. \\
CNE-U & Seccion 050, pagina 69 & Apoya la revision de circuitos derivados en unidades de vivienda. \\
\bottomrule
\end{longtable}

Las evidencias normativas estan guardadas en la carpeta de normativa extraida del avance. Si alguna pagina no permite extraccion de texto, se revisara visualmente como imagen.

\section*{6. Supuestos de potencia para el primer avance}

\begin{longtable}{L{7cm}R{4cm}L{3cm}}
\toprule
\textbf{Elemento} & \textbf{Potencia referencial} & \textbf{Estado} \\
\midrule
\endfirsthead
\toprule
\textbf{Elemento} & \textbf{Potencia referencial} & \textbf{Estado} \\
\midrule
\endhead
Foco LED & 12 W & Preliminar \\
Tomacorriente general & 180 W & Preliminar \\
Tomacorriente de cocina o servicio & 300 W & Preliminar \\
Microondas / horno electrico & 1200 W & Referencial \\
Lavadora & 800 W & Referencial \\
Waflera & 1000 W & Referencial \\
Bomba de agua & 750 W & Referencial \\
\bottomrule
\end{longtable}

\section*{7. Levantamiento preliminar de cargas}

\begin{longtable}{L{4.2cm}L{4.2cm}R{1.6cm}R{2.2cm}R{2.2cm}}
\toprule
\textbf{Ambiente / zona} & \textbf{Carga} & \textbf{Cant.} & \textbf{W unit.} & \textbf{W total} \\
\midrule
\endfirsthead
\toprule
\textbf{Ambiente / zona} & \textbf{Carga} & \textbf{Cant.} & \textbf{W unit.} & \textbf{W total} \\
\midrule
\endhead
Primer piso & Focos LED & 10 & 12 & 120 \\
Cuartos primer piso & Tomacorrientes generales & 6 & 180 & 1080 \\
Cocina primer piso & Tomacorriente cocina & 1 & 300 & 300 \\
Segundo piso & Focos LED & 11 & 12 & 132 \\
Sala, cuartos y bano segundo piso & Tomacorrientes generales asumidos & 13 & 180 & 2340 \\
Cocina amplia segundo piso & Tomacorrientes cocina/servicio & 4 & 300 & 1200 \\
Cocina amplia segundo piso & Microondas / horno electrico & 1 & 1200 & 1200 \\
Cocina amplia segundo piso & Lavadora & 1 & 800 & 800 \\
Cocina amplia segundo piso & Waflera & 1 & 1000 & 1000 \\
Exterior primer piso & Bomba de agua & 1 & 750 & 750 \\
\midrule
\textbf{Total preliminar} &  &  &  & \textbf{8922} \\
\bottomrule
\end{longtable}

La tabla anterior es solo una base de trabajo. Los tomacorrientes generales del segundo piso se han asumido para poder preparar el calculo inicial; se ajustaran cuando se defina el croquis.

\section*{8. Circuitos preliminares recomendados}

\begin{longtable}{L{1.2cm}L{3.7cm}L{7.2cm}L{1.7cm}}
\toprule
\textbf{Cto.} & \textbf{Uso} & \textbf{Ambientes / cargas} & \textbf{Prior.} \\
\midrule
\endfirsthead
\toprule
\textbf{Cto.} & \textbf{Uso} & \textbf{Ambientes / cargas} & \textbf{Prior.} \\
\midrule
\endhead
C1 & Alumbrado primer piso & Cuartos, cocina pequena, pasadizo y escalera/plataforma & Alta \\
C2 & Tomacorrientes primer piso & Cuarto 1 y cuarto 2 & Alta \\
C3 & Cocina primer piso & Tomacorriente para licuadora y artefactos pequenos & Alta \\
C4 & Alumbrado segundo piso & Sala, cuartos, cocina amplia y bano & Alta \\
C5 & Tomacorrientes generales segundo piso & Sala, cuartos, cuarto de uso varios y bano & Alta \\
C6 & Cocina amplia segundo piso & Tomacorrientes de cocina, microondas/horno y waflera & Alta \\
C7 & Lavadora segundo piso & Carga de servicio dedicada & Media \\
C8 & Bomba de agua exterior & Bomba del pozo & Alta \\
\bottomrule
\end{longtable}

\section*{9. Resumen preliminar de demanda}

Para un primer avance se agrupan las cargas por tipo y se aplican factores academicos preliminares. Estos factores deben corregirse si el docente exige usar una tabla especifica del CNE-U.

\begin{longtable}{L{4.5cm}R{3cm}R{2.8cm}R{3cm}}
\toprule
\textbf{Tipo de carga} & \textbf{Potencia W} & \textbf{Factor} & \textbf{Demanda W} \\
\midrule
\endfirsthead
\toprule
\textbf{Tipo de carga} & \textbf{Potencia W} & \textbf{Factor} & \textbf{Demanda W} \\
\midrule
\endhead
Alumbrado & 252 & 1.00 & 252 \\
Tomacorrientes generales & 3420 & 0.70 & 2394 \\
Cocina y servicio & 4500 & 0.80 & 3600 \\
Bomba de agua & 750 & 1.00 & 750 \\
\midrule
\textbf{Total} & \textbf{8922} &  & \textbf{6996} \\
\bottomrule
\end{longtable}

\[
I = \frac{P}{V \times fp} = \frac{6996}{220 \times 0.90} \approx 35.33\ A
\]

Por el resultado preliminar, el interruptor general podria evaluarse en el orden de 40 A o 50 A, sujeto a correccion por potencia real de equipos, longitud de alimentador, criterio del docente y tabla normativa aplicable.

\section*{10. Pendientes para cerrar el Capitulo I}

\begin{itemize}
  \item Confirmar nombre del docente.
  \item Confirmar provincia y distrito exactos.
  \item Dibujar croquis del primer y segundo piso.
  \item Estimar dimensiones aproximadas de ambientes.
  \item Definir ubicacion del tablero general.
  \item Confirmar cantidad final de tomacorrientes en sala, cuartos y cuarto de uso varios.
  \item Confirmar potencias reales o comerciales de bomba, lavadora, microondas/horno y waflera.
  \item Decidir si la canalizacion sera empotrada, superficial o mixta.
\end{itemize}

\end{document}
